\chapter{整台机器}
% 完整版：chapters/10_synthesis.tex

\pencilfig{10}{\pencilart{10}}{整台机器：数据总线，其中的刹车，以及上方的四个控制广播频道。}

我们每次拿出一种神经元，问它给机器添了什么。现在该把零件装到一起了。

装出来的是一台三层的机器。最底层是由\nt{GLUT}神经元组成的巨大电话网：数据在上面传送。旁边是抑制性的\nt{GABA}神经元：它们管理信号流，决定什么通过、何时通过、以多大音量通过。而在这一切之上，是几座电台，每座都有自己的旋钮。\nt{DOPA}说明哪些连接的改变应当保留。\nt{SERO}维持全局水平，按照假说还掌管耐心的时间跨度。\nt{NORA}在探索和决定之间做选择。\nt{ACET}——在写入和读取之间。每座电台实际上是一小束频道，而不是一根导线，但相对于八百六十亿个神经元，频道仍然只有一小撮。

\section*{一个故事}

让我们跟随一个事件走遍整台机器。动物早就知道：按下杠杆 A，就能得到果汁。有一天世界变了：A 不再管用，带来果汁的是杠杆 B，而动物几乎没用过它。

按下 A 后果汁没来——多巴胺报以低谷，选择 A 的那些连接减弱。低谷一再出现，误差比平常大——按照假说，去甲肾上腺素报告：世界变了。对比度旋钮调低，选择变得更柔和，噪声更常推动动物去试 B。乙酰胆碱把网络切换到写入模式。试一下 B，带来了意外的果汁——多巴胺爆发，选择 B 的连接增强。几次尝试之后，预期追上了新的世界，意外没有了，各个旋钮回到原位。

注意：没有哪个旋钮知道其他旋钮在做什么。每个旋钮只对自己的标志作出反应，动作的顺序是自然形成的，就像在异步电路里，每个模块在输入就绪时触发。每个旋钮各自这样表现，这大部分已经测量到。它们能协调一致地共同工作——目前还是假说。

\section*{它和你的笔记本电脑哪里不同}

笔记本电脑有统一的时钟节拍——大脑没有。笔记本电脑的元件是可靠的——突触会丢失一半以上的信号。笔记本电脑的存储器与处理器是分开的——大脑的记忆就在同样的那些连接里。笔记本电脑（如果它学习的话）靠误差反向传播来学习——大脑靠局部规则和一个广播式的“比预期的好”。最重要的是：笔记本电脑光一个处理器就要耗掉几十瓦，而大脑全部加起来只要二十瓦。

\section*{我们还不知道什么}

我们不知道皮层在多大程度上利用了咔嗒的精确时刻。不知道在所有人只有一位老师的情况下，它是如何调好多层回路中数十亿个连接的。还不完全明白电台们传送的是什么。不知道在没有统一节拍的情况下，大脑相距遥远的各个部分怎样协调工作。而且至今没有人能造出一台只用二十瓦就做到同样事情的机器。

接下来，本书将走出这个核心：去看机器的输入和输出，看怎样调试它，看睡眠，以及芯片上由活神经元构成的机器。

\fullversion{10}
